% Hikaru: implementation. Kelvin: full cluster run and analysis.
\subsection{\expt{a102}: Lookup indexes}
\label{sec:a102}
The first experiment added three indexes:
\begin{lstlisting}
CREATE INDEX idx_orders_customer_recent
    ON orders (merchant_id, customer_id, created_at DESC);
CREATE INDEX idx_items_order ON order_items (order_id);
CREATE INDEX idx_inventory_product ON inventory (product_id);
\end{lstlisting}
The orders index follows T3's merchant/customer filtering and recent-order
access pattern. The others support order-item and inventory lookups. This
ordering follows PostgreSQL's guidance on leading equality columns in
multicolumn B-tree indexes \citep{pg-multicolumn}.

Kelvin ran a102 as job 965657. Its throughput was
3.250/3.453/3.543/3.320/3.782/3.398 TPS at 1--64 clients. All measured phases
exited normally and recorded no serialization failures. However, the 16-client
warmup aborted, excluding that case from comparisons requiring both phases
to complete.

Against the preserved baseline, eligible concurrent cases appeared 49--57\%
faster. Against the fresh baseline, differences were only approximately
$-3\%$ to $+6\%$. Different node allocations and single observations prevent
a large overall throughput claim attributable to the indexes.

At one client, T7's mean latency decreased from 51.114 to 13.709\,ms, while
T3 and T4 remained expensive at 2,668.472 and 454.362\,ms. The indexes helped
particular operations but did not remove the dominant distributed-read cost.
\evidence{\expt{a102}, job 965657; preserved a0 job 918524 and fresh a0 job
977740. Source identity was checked against the saved runtime hashes.}
